% Corrigé intégré automatiquement pour la banque Exercices.
% Source : C2_MettreEnOeuvreDemarche/C2_07_PFS/56_RobotAvion/56_RobotAvion.tex
\begin{corrigebox}[Corrigé]
\CorrigeQuestion{1}
\begin{center}
\includegraphics[width=.5\linewidth]{56_01_cor}
\end{center}
\CorrigeQuestion{2}
En isolant l'ensemble $\Sigma = \{2+3+4+E_1\}$ on ne fera apparaître \textbf{QUE} $C_{12}$ et les actions mécaniques extérieures. Les actions de liaison 2--3, 3--4, 4--E1 \textbf{n'interviendont pas}.
\CorrigeQuestion{3}
\begin{center}
\includegraphics[width=.5\linewidth]{56_02_cor}
\end{center}
Bilan des actions mécaniques :
\begin{itemize}
\item pivot 1--2  et couple moteur de 1 sur 2;

\item pesanteur sur 2 : $\torseurstat{T}{\text{pes}}{2}$ $=\torseurl{-M_2g \vect{y_1} }{\vect{0}}{G_2}$.
On a par ailleurs $\vectm{O_2}{\text{pes}}{2} = \vect{O_2 G_2} \wedge-M_2g \vect{y_1}$
$= \dfrac{1}{2}L_3\vect{x_2} \wedge-M_2g \vect{y_1}$
$= -\dfrac{1}{2}M_2gL_3 \cos \theta_{12} \vect{z_0} $;

\item pesanteur sur 3 et 4 : $\torseurstat{T}{\text{pes}}{3+4}$ $=\torseurl{-M_{34} g \vect{y_1} }{\vect{0}}{G_3}$.
On a par ailleurs $\vectm{O_2}{\text{pes}}{3+4} = \vect{O_2 G_3} \wedge-M_{34}g \vect{y_1}$
$= \left( L_3\vect{x_2} +\dfrac{1}{3}L_4\vect{x_3}+L_5\vect{y_3}\right) \wedge-M_{34}g \vect{y_1}$
$= -M_{34}g\left( L_3\vect{x_2}\wedge \vect{y_1} +\dfrac{1}{3}L_4\vect{x_3}\wedge \vect{y_1}+L_5\vect{y_3}\wedge \vect{y_1} \right) $
$= -M_{34}g\left( L_3\cos\theta_{12} +\dfrac{1}{3}L_4\cos\theta_{12}-L_5\sin\theta_{13}\right) \vect{z_0}$;

\item pesanteur sur $E_1$ : $\torseurstat{T}{\text{pes}}{E1}$ $=\torseurl{-M_{E1} g \vect{y_1} }{\vect{0}}{G_5}$.
On a par ailleurs $\vectm{O_2}{\text{pes}}{E1} = \vect{O_2 G_5} \wedge-M_{E1}g \vect{y_1}$;
$= -M_{E1}g\left( L_3 \vect{x_2}+L_4\vect{x_3}+L_5 \vect{x_3}+l_5 \vect{y_3}+L_7 \vect{x_5} \right)\wedge \vect{y_1}$
$= -M_{E1}g\left( L_3 \cos\theta_{12}+L_4\cos\theta_{13}+L_5 \cos\theta_{13}-l_5 \sin\theta_{13}+L_7 \cos\theta_{15}\right) \vect{z_0}$;

\item effort presseur + perçage $\torseurstat{T}{\text{Tronçon }}{E_1}=\torseurcol{-F-P}{0}{0}{0}{0}{0}{P, \rep{5}}$.
On a par ailleurs $\vectm{O_2}{\text{Tronçon}}{E1} = \vect{O_2 P} \wedge(-F-P) \vect{x_5}$
$= \left(  L_3 \vect{x_2}+L_4\vect{x_3}+L_5 \vect{x_3}+l_5 \vect{y_3}+L_8 \vect{x_5} \right)\wedge(-F-P) \vect{x_5}$
$= -(F+P)\left(  L_3 \vect{x_2}\wedge \vect{x_5}+L_4\vect{x_3}\wedge \vect{x_5}+L_5 \vect{x_3}\wedge \vect{x_5}+l_5 \vect{y_3}\wedge \vect{x_5} \right)\wedge $
$= -(F+P)\left(  L_3 \sin\left(\theta_{15}-\theta_{12} \right)+ \left(L_4+L_5\right)\sin\left(\theta_{15}-\theta_{13} \right)+l_5 \sin\left(\theta_{15}-\theta_{13} -\dfrac{\pi}{2} \right)  \right)\vect{z_0}$.
\end{itemize}
\CorrigeQuestion{4}
Pour ne pas faire apparaître les actions de la liaison 1--2 il faudra réaliser un théorème du moment statique en $O_2$ en projection sur $\vect{z_2}$ (la liaison pivot n'a pas de composante en ce point et sur cette projection).

On a donc : 
$
 -\dfrac{1}{2}M_2gL_3 \cos \theta_{12}
  -M_{34}g\left( L_3\cos\theta_{12} +\dfrac{1}{3}L_4\cos\theta_{12}-L_5\sin\theta_{13}\right)
  -M_{E1}g\left( L_3 \cos\theta_{12}+L_4\cos\theta_{13}+L_5 \cos\theta_{13}-l_5 \sin\theta_{13}+L_7 \cos\theta_{15}\right) 
  -(F+P)\left(  L_3 \sin\left(\theta_{15}-\theta_{12} \right)+ \left(L_4+L_5\right)\sin\left(\theta_{15}-\theta_{13} \right)+l_5 \sin\left(\theta_{15}-\theta_{13} -\dfrac{\pi}{2} \right)  \right) 
  = 0
$.
\CorrigeQuestion{5}
On a 
$
C_{12} -\dfrac{1}{2}M_2gL_3 \cos \theta_{12}
  -M_{34}g\left( L_3\cos\theta_{12} +\dfrac{1}{3}L_4\cos\theta_{12}\right)
  -M_{E1}g\left( L_3 \cos\theta_{12}+L_4+L_5+L_7 \cos\theta_{15}\right) 
  -(F+P)\left(  L_3 \sin\left(\theta_{15}-\theta_{12} \right)+ \left(L_4+L_5\right)\sin\left(\theta_{15} \right)-l_5 \cos\left(\theta_{15}  \right)  \right) 
  = 0
$.
\CorrigeQuestion{6}
On a 
$
C_{12} -\dfrac{1}{2}M_2gL_3 \cos \theta_{12}
  -M_{34}g\left( L_3\cos\theta_{12} +\dfrac{1}{3}L_4\cos\theta_{12}\right)
  -M_{E1}g\left( L_3 \cos\theta_{12}+L_4+L_5 \right) 
  -(F+P)\left(  L_3 \cos\left(\theta_{12} \right)- \left(L_4+L_5\right)  \right) 
  = 0
$.
Soit 
$
C_{12} -\dfrac{1}{4}M_2gL_3 
  -M_{34}g\dfrac{1}{2}\left( L_3 +\dfrac{1}{3}L_4\right)
  -M_{E1}g\left( L_3 \dfrac{1}{2}+L_4+L_5 \right) 
  -(F+P)\left(  L_3 \dfrac{1}{2}- \left(L_4+L_5\right)  \right) 
  = 0
$.
Au final 
$
C_{12} 
= \dfrac{1}{4}M_2gL_3 
  +M_{34}g\dfrac{1}{2}\left( L_3 +\dfrac{1}{3}L_4\right)
  +M_{E1}g\left( L_3 \dfrac{1}{2}+L_4+L_5 \right) 
  +(F+P)\left(  L_3 \dfrac{1}{2}- \left(L_4+L_5\right)  \right) 
$.
\CorrigeQuestion{7}
$C_{12} = \SI{6230}{Nm}$. Compatible avec le cahier des charges.
\end{corrigebox}
